\subsection{两个对应的复合}

让 \(G, G'\) 是两个图。设 A 表示集合 \(pr_1G\) 并设 C 表示集合 \(pr_2G'\) 。关系 \((\exists y)((x, y) \in G\) 并且 \((y, z) \in G')\) 蕴含 \((x, z) \in A \times C\) ，因此关于 x 和 z 有一个图。

定义6. 设G，G'为两个图。关系（关于x和z）的图 \((\exists y)((x, y) \in G \text{ 且 } (y, z) \in G')\) 称为G'与G的复合，记作G' o G（有时也记作G'G）。

命题3. 设G、G'为两个图。则G' o G的逆为 \(\overline{G} \circ \overline{G'}\) .

对于关系“ \((x, y) \in G\) 并且 \((y, z) \in G'\) ”等价于

\[
(z, y) \in \overline {{\mathbf {G} ^ {\prime}}} \text {   且   } (y, x) \in \overline {{\mathbf {G}}}
\]

命题4。设 \(G_{1}\) , \(G_{2}\) , \(G_{3}\) 是图。则

\[
\left(\mathrm{G} _ {3} \circ \mathrm{G} _ {2}\right) \circ \mathrm{G} _ {1} = \mathrm{G} _ {3} \circ \left(\mathrm{G} _ {2} \circ \mathrm{G} _ {1}\right).
\]

关系 \((x, t) \in (\mathbf{G}_3 \circ \mathbf{G}_2) \circ \mathbf{G}_1\) 等价于关系

\[
(\exists y) ((x, y) \in \mathrm{G} _ {1} \text {   且   } \exists z ((y, z) \in \mathrm{G} _ {2} \text {   且   } (z, t) \in \mathrm{G} _ {3}))
\]

因此（根据C33（第一章，§4，第3节））与关系

\[
(\exists y) (\exists z) ((x, y) \in \mathrm{G} _ {1} \text {   且   } (y, z) \in \mathrm{G} _ {2} \text {   且   } (z, t) \in \mathrm{G} _ {3}).\tag{1}
\]

同样地，关系 \((x, t) \in \mathbf{G}_3 \circ (\mathbf{G}_2 \circ \mathbf{G}_1)\) 等价于

\[
(\exists z) (\exists y) ((x, y) \in \mathrm{G} _ {1} \text {   且   } (y, z) \in \mathrm{G} _ {2} \text {   且   } (z, t) \in \mathrm{G} _ {3}).\tag{2}
\]

但关系(1)与(2)是等价的；因此结论成立。

¶ 图 \(G_{3} \circ (G_{2} \circ G_{1})\) 处记作 \(G_{3} \circ G_{2} \circ G_{1}\) . 类似地，如果 \(G_{1}, G_{2}, G_{3}, G_{4}\) 是图，我们令

\[
\mathrm{G} _ {4} \circ (\mathrm{G} _ {3} \circ \mathrm{G} _ {2} \circ \mathrm{G} _ {1}) = \mathrm{G} _ {4} \circ \mathrm{G} _ {3} \circ \mathrm{G} _ {2} \circ \mathrm{G} _ {1};
\]

等等。

命题5. 设 G, G' 为图，且 A 为一个集合。则

\[
(\mathbf {G} ^ {\prime} \circ \mathbf {G}) \langle \mathbf {A} \rangle = \mathbf {G} ^ {\prime} \langle \mathbf {G} \langle \mathbf {A} \rangle \rangle .
\]

因为根据 C33（第一章，第4节，第3号）的关系 \(z \in (\mathbf{G}' \circ \mathbf{G})\langle \mathbf{A} \rangle\) 等价于

\[
(\exists y) ((\exists x) (x \in A \text {   且   } (x, y) \in G) \text {   且   } (y, z) \in G ^ {\prime})
\]

并且因此等价于 \((\exists y)(y \in G \langle A \rangle\) 并且 \((y, z) \in G')\) ；因此结论成立。

如果 G 和 \(G'\) 是两个图，我们有 \(\mathrm{pr}_{1}(\mathrm{G}' \circ \mathrm{G}) = \mathrm{G}\langle\mathrm{pr}_{1}\mathrm{G}'\rangle\) ，以及 \(\mathrm{pr}_{2}(\mathrm{G}' \circ \mathrm{G}) = \mathrm{G}'\langle\mathrm{pr}_{2}\mathrm{G}\rangle\) . 例如，要证明这些关系中的第二个，只需注意到关系 \(z \in \mathrm{pr}_2(G' \circ G)\) 等价于 \((\exists x)((x, z) \in G' \circ G)\) 并且因此

\[
(\exists y) ((\exists x) ((x, y) \in G) \quad \text { 且 } \quad (y, z) \in G ^ {\prime});
\]

但这等价于 \(z \in G' \langle pr_{2}G \rangle\) .

如果 G 是一个图，且 X 是一个集合，使得 \(X \subset pr_{1}G\) ，我们有

\[
\mathbf {X} \subset \overline {{\mathbf {G}}} \langle \mathbf {G} \langle \mathbf {X} \rangle \rangle .
\]

由假设，关系 \(x \in \mathbf{X}\) 蕴含 \((\exists y)((x, y) \in \mathbf{G})\) ；而 \((x, y) \in \mathbf{G}\) 等价于 \((y, x) \in \overline{\mathbf{G}}\) ，另一方面， \((x, y) \in \mathbf{G}\) 蕴含 \((\exists z)(z \in \mathbf{X} \text{ 且 } (z, y) \in \mathbf{G})\) ；因此 \(x \in \mathbf{X}\) 蕴含

\[
(\exists y) ((\exists z) (z \in X \text { 且 } (z, y) \in G) \text { 且 } (y, x) \in \overline {{G}} ^ {- 1}),
\]

也就是说， \(x \in \overline{\mathbf{G}}\langle \mathbf{G}\langle \mathbf{X}\rangle\rangle\) .

显然，如果 \(G_{1}\) , \(G_{2}\) , \(G_{1}^{\prime}\) , \(G_{2}^{\prime}\) 是图，则关系 \(G_{1} \subset G_{2}\) 并且 \(G_{1}^{\prime} \subset G_{2}^{\prime}\) 蕴含 \(G_{1}^{\prime} \circ G_{1} \subset G_{2}^{\prime} \circ G_{2}\) .

让 \(\Gamma = (G, A, B)\) 并且 \(\Gamma' = (G', B, C)\) 是两个对应，使得 \(\Gamma\) 的目标与 \(\Gamma'\) 的源相同。由上述讨论我们有 \(\mathrm{pr}_1(G' \circ G) \subset \mathrm{pr}_1G \subset A\) 并且 \(\mathrm{pr}_2(G' \circ G) \subset \mathrm{pr}_2G' \subset C\) ；因此我们可以给出如下定义：

定义7. 设 \(\Gamma = (G, A, B)\) 和 \(\Gamma' = (G', B, C)\) 是两个对应，使得 \(\Gamma\) 的目标与 \(\Gamma'\) 的源相同。则对应 \((G' \circ G, A, C)\) 称为 \(\Gamma'\) 与 \(\Gamma\) 的复合，并记作 \(\Gamma' \circ \Gamma\) （有时也记作 \(\Gamma' \Gamma\) ).

由命题5立即得出，若X是A的子集，则我们有 \((\Gamma'\circ\Gamma)\langle X\rangle=\Gamma'\langle\Gamma\langle X\rangle\rangle\) 。此外，由于 \(\overline{\Gamma}^{1}\) 的目标与 \(\overline{\Gamma}\) 的源相同， \(\Gamma'\circ\Gamma\) 的逆为 \(\overline{\Gamma}\circ\overline{\Gamma}'\) （根据命题3）。

定义8。若A是一个集合，则该集合 \(\Delta_{\mathbf{A}}\) 所有形如 \((x, x)\) ，其中 \(x \in \mathbf{A}\) ，称为 \(\mathbf{A} \times \mathbf{A}\) .

的对角线。显然我们有 \(\mathrm{pr}_1\Delta_{\mathbf{A}} = \mathrm{pr}_2\Delta_{\mathbf{A}} = \mathbf{A}\) 该对应

\[
\mathbf {I} _ {\mathbf {A}} = (\Delta_ {\mathbf {A}}, \mathbf {A}, \mathbf {A})
\]

称为 A 上的恒等对应；它是自身的逆。

如果 \(\Gamma\) 是 A 与 B 之间的一个对应，并且若 \(I_{A}, I_{B}\) 分别是 A、B 上的恒等对应，则 \(\Gamma \circ I_{A} = I_{B} \circ \Gamma = \Gamma\) .

